% Nouveau EGA V (cote 25): archive cote 25, site page 115 (PDF page 116).
% typescript. Draft compared: none.
% AI-transcribed by Claude Opus 5.5 (claude-opus-5-5, adaptive thinking, effort high), Message Batches API, batch ega25b (msgbatch_01A7n5YuXubUPNK1zsepi8RG), 2026-11-13; checked against both supplied scan parts; see CHECK_DECISIONS.tsv.
\scanpage{25/115}
\mspage{18}
changement de base en S; d'autre part, X-Z \var{\inserted[typed]{(supposé ouvert, p ex T/S plat lpf.}}{(supposé ouvert, p ex T/S plat lpf.}{insertion}{ega25_25_115-v0001} est universellement
schématiquement dense dans X rel à S\var{\inserted{lorsque $T \to S$ est surjectif\unsure{)}}}{ lorsque $T \to S$ est surjectif\unsure{)}}{insertion}{ega25_25_115-v0002}. Ce dernier fait résulte \unsure{du|de} \var{\xout{la compa}}{}{typewriter strike-out}{ega25_25_115-v0003}%
\var{\xout{tibilité avec changement de base}}{}{typewriter strike-out}{ega25_25_115-v0004} par. 11 et du fait que pour tout $s \hand{\in} S$,
$Z_s$ ne contient aucun point de $X_s$ associé à $\underline{O}_{X_s}$ (en effet, le support d'un
diviseur \var{\deleted{de Cartier}}{}{deletion}{ega25_25_115-v0005} sur $X_s$ ne contient aucun tel point).

Dans le cas où T \var{\deleted{devient}}{}{deletion}{ega25_25_115-v0006}\var{\inserted{est}}{est}{insertion}{ega25_25_115-v0007} un fibré projectif $Q=P(F)$, $\mathrm{Div}_{Q/S}$
étant représentable par la somme des $P(\mathrm{Sym}^n(\overset{\hand{\vee}}{F})) = P(n)$, on trouve
un morphisme $X-Z \hand{\longrightarrow} P(n)$ . On dit\var{\xout{\unsure{a|q}}}{}{typewriter strike-out}{ega25_25_115-v0008} que la famille de diviseurs
D est \uline{de degré} n, si le morphisme précédent se factorise par P(n); si
$X \hand{\neq} \hand{\emptyset}$ , d'où $X-Z \hand{\neq} \hand{\emptyset}$ , le n en question est bien déterminé par D. D'ail%
leurs pour définir cette notion, on n'a pas besoin à la rigueur du résul-%
tat de représentabilité énoncé\var{\xout{s}}{}{typewriter strike-out}{ega25_25_115-v0009} plus haut, mais seulement d'avoir défini
les monomorphismes canoniques
\[ P(n) \hand{\longrightarrow} \mathrm{Div}_{\overset{\hand{\vee}}{P}/S} \]
(NB on pose $P = P(\overset{\hand{\vee}}{F})$, d'où $Q = P(F) = \overset{\hand{\vee}}{P}$ avec les notations des n°s précé-%
dents). \var{\deleted{Cela pourra être fait à titre d'exemple dans le n° préconisé au
par. 19}}{}{deletion}{ega25_25_115-v0010}. On appelle \uline{système linéaire de diviseurs sur} X/S , \uline{paramétré
par le fibré projectif} $Q = \overset{\hand{\vee}}{P}$ , \uline{toute famille de diviseurs} \uline{sur} X/S
\uline{paramétrée par} Q, \uline{qui soit de degré} 1, i.e. \uline{définisse}
\[ \hand{f} : X-Z \hand{\longrightarrow} P \qquad . \]
Ainsi, à un tel système \var{\deleted{\var{\xout{\unsure{li}}}{}{typewriter strike-out}{ega25_25_115-v0012}}}{}{deletion}{ega25_25_115-v0011} linéaire de diviseurs\var{\inserted{(et lorsque les fibres de $\check{P}$ sont tous $\neq \emptyset$)}}{(et lorsque les fibres de $\check{P}$ sont tous $\neq \emptyset$)}{insertion}{ega25_25_115-v0013}, est associé canoniquement
une application rationelle de X dans un fibré projectif,\var{\inserted{savoir un ``\var{\deleted{\illegible{1}}}{}{deletion}{ega25_25_115-v0015} pseudo-morphisme rel.\unsure{:|à} S''}}{savoir un ``\var{\deleted{\illegible{1}}}{}{deletion}{ega25_25_115-v0015} pseudo-morphisme rel.\unsure{:|à} S''}{insertion}{ega25_25_115-v0014} \hand{(}en fait, un peu
mieux même qu'une application rationelle\hand{)}. Par construction même,
$D\,\hand{|}(X-Z)x_S\overset{\hand{\vee}}{P}$ n'est autre que l'image inverse par $f x id_{\overset{\hand{\vee}}{P}}$ du diviseur
canonique (le diviseur d'in\unsure{c}id\var{\deleted{\unsure{e}}}{}{deletion}{ega25_25_115-v0016}ence )\hand{\uline{H}} sur $Px_S\overset{\hand{\vee}}{P}$ , Donc la connaissance
de $f: X-Z \hand{\longrightarrow} P$ permet de réconstituer \var{\xout{la fa}}{}{typewriter strike-out}{ega25_25_115-v0017} tout au moins la famil-%
le de diviseurs de X-Z induite par D , - donc si la famille est sans
\mspage{115}
